\section[Rung two: signed one-hop relaxation]
{Rung two: signed one-hop relaxation pays
$\Theta(1/(\sqrt\alpha\,\epsppr))$}
\label{sec:ladder-signed-lower-bound}

Dropping $r\ge0$ destroys the order argument of
\Cref{sec:ladder-monotone-lower-bound}, since $\Psi\ge0$ was exactly
\ref{it:ladder-f2} applied to a nonnegative residual. It is replaced here by a
quadratic energy argument, which is why the class is restricted to relaxations
$\delta=\omega r_{u}$ with $\omega\in(0,2)$: by
\Cref{lem:ladder-energy} that interval \emph{is} the dissipativity condition.

\begin{lemma}[Energy identity; \statusproveddraft]
\label{lem:ladder-energy}
Put $\Phi(e):=\tfrac12\,e^{\trans}D^{-1}He$; the matrix
$D^{-1}H=D^{-1/2}\Hsym D^{-1/2}$ is symmetric positive definite, so $\Phi$ is
a norm on the error. A $\push(u,\delta)$ changes it by
\begin{equation}
    \Phi(e)-\Phi(e^{+})=\frac{\delta}{d_{u}}\Bigl(r_{u}-\frac{\delta}{2}\Bigr),
    \label{eq:ladder-energy-identity}
\end{equation}
which for $\delta=\omega r_{u}$ equals
$\tfrac{\omega(2-\omega)}{2}\,r_{u}^{2}/d_{u}\ge0$. Hence $\Phi$ is
non-increasing along every $\Rpm$ path, and $\omega\in(0,2)$ is exactly the
dissipativity condition. Initially
$\Phi_{0}=\galpha\pi_{v}/(2d_{v})$, which on the star is $\alpha/(2m)$.
\end{lemma}

\begin{proof}
Write $e^{+}=e-\delta\,\eunit{u}$. Since $(D^{-1}He)_{u}=r_{u}/d_{u}$ and
$(D^{-1}H)_{uu}=1/d_{u}$, expanding the quadratic gives
\eqref{eq:ladder-energy-identity}. For $\Phi_{0}$ use $e^{0}=\pi$ and
$H\pi=\galpha e_{v}$, so
$\Phi_{0}=\tfrac12\pi^{\trans}D^{-1}\galpha e_{v}=\galpha\pi_{v}/(2d_{v})$;
on the star \eqref{eq:ladder-star-closed-forms} gives $\alpha/(2m)$.
\end{proof}

\begin{lemma}[Fast-mode cap; \statusproveddraft]
\label{lem:ladder-fast-mode}
Let $y:=D^{-1/2}e$, so that $\Phi=\tfrac12y^{\trans}\Hsym y$ with $\Hsym$ as
in \eqref{eq:ladder-Hsym}. For any unit eigenvector $q$ of $\Hsym$ with
eigenvalue $\mu$, the coefficient $\widehat y_{q}:=\inner{y}{q}$ satisfies, at
every time,
\begin{equation}
    \tfrac12\,\mu\,\widehat y_{q}^{2}\ \le\ \Phi\ \le\ \Phi_{0},
    \qquad\text{i.e.}\qquad
    \abs{\widehat y_{q}}\ \le\ \sqrt{2\Phi_{0}/\mu}.
    \label{eq:ladder-fast-mode-cap}
\end{equation}
\end{lemma}

\begin{proof}
$\Phi=\tfrac12\sum_{k}\mu_{k}\widehat y_{k}^{2}$ with every $\mu_{k}>0$ by
\eqref{eq:ladder-Hsym}; drop all terms but one and apply
\Cref{lem:ladder-energy}.
\end{proof}

\paragraph{The star's exact two-dimensional reduction.}
On $K_{1,m}$ the matrix $D^{-1/2}AD^{-1/2}$ has the symmetric eigenpairs
$\lambda=+1$ and $\lambda=-1$ with unit eigenvectors
$q_{\pm}=(\sqrt m,\pm1,\ldots,\pm1)^{\trans}/\sqrt{2m}$, together with $m-1$
leaf-difference eigenvectors of eigenvalue $0$ that \emph{vanish at the
centre}. Hence $\Hsym$ has the slow mode $q_{+}$ with
$\mu_{-}=1-\calpha=\galpha$, the fast mode $q_{-}$ with
$\mu_{+}=1+\calpha=2/(1+\alpha)$, and an inert middle block. Seed-symmetric
dynamics therefore live in the two-dimensional $(e_{c},E_{L})$ plane, where
$E_{L}:=\sum_{l\in L}e_{l}$ and
$\widehat y_{\pm}=(e_{c}\pm E_{L})/\sqrt{2m}$ up to sign convention.

\begin{lemma}[Equal-magnitude displacement cap; \statusproveddraft]
\label{lem:ladder-displacement}
A centre relaxation $\push(c,\delta)$ changes $y$ by
$-(\delta/\sqrt m)\,\eunit{c}$, hence moves the two modes by equal magnitudes,
\begin{equation}
    \abs{\Delta\widehat y_{+}}=\abs{\Delta\widehat y_{-}}
    =\frac{\abs{\delta}}{\sqrt{2m}} .
    \label{eq:ladder-equal-magnitude}
\end{equation}
The slow coefficient carries the answer, but the fast coefficient is confined
to the interval of \Cref{lem:ladder-fast-mode}; applying that lemma before and
after the operation gives, at every centre relaxation, for every
$\omega\in(0,2)$, pathwise,
\begin{equation}
    \abs{\delta}\ \le\ \sqrt{2m}\cdot2\sqrt{2\Phi_{0}/\mu_{+}}
    \ =\ 2\sqrt{\alpha(1+\alpha)} .
    \label{eq:ladder-displacement-cap}
\end{equation}
\end{lemma}

\begin{proof}
$\Delta e=-\delta\eunit{c}$ and $y=D^{-1/2}e$ give
$\Delta y=-(\delta/\sqrt{d_{c}})\eunit{c}$ with $d_{c}=m$; pairing with
$q_{\pm}$, whose centre coordinate is $\sqrt m/\sqrt{2m}$ in absolute value,
gives \eqref{eq:ladder-equal-magnitude}. Since $\widehat y_{-}$ lies in
$[-\sqrt{2\Phi_{0}/\mu_{+}},\,\sqrt{2\Phi_{0}/\mu_{+}}]$ both before and after
the operation, $\abs{\Delta\widehat y_{-}}\le2\sqrt{2\Phi_{0}/\mu_{+}}$;
multiply by $\sqrt{2m}$ and substitute $\Phi_{0}=\alpha/(2m)$,
$\mu_{+}=2/(1+\alpha)$.
\end{proof}

\begin{remark}[Bursts do not evade the cap]
\label{rem:ladder-bursts}
Consecutive operations at the same vertex compose to a single relaxation with
$\omega_{\mathrm{eff}}=1-\prod_{j}(1-\omega_{j})\in(0,2)$, so a burst at the
centre is itself a member operation and inherits
\eqref{eq:ladder-displacement-cap}. The fast mode is moreover
$1-\bigO(\alpha)$-saturated already at $t=0$, so the cap is not merely
asymptotic.
\end{remark}

\begin{theorem}[Class lower bound for $\Rpm$; \statusproveddraft]
\label{thm:ladder-signed-lower-bound}
Take the centre-seeded star \eqref{eq:ladder-hard-instance} with
$\epsppr\le1/16$ and $\alpha\in(0,1)$. Every member of $\Rpm$ --- arbitrary
adaptive, randomized or per-vertex-$\omega$ policy --- that stops with
$\errop(z)\le\epsppr$ performs
\begin{equation}
    N_{c}\ \ge\ \frac{\pi_{c}-\epsppr m}{2\sqrt{\alpha(1+\alpha)}}
    \ \ge\ \frac{3}{16\sqrt{\alpha(1+\alpha)}}
    \label{eq:ladder-signed-count}
\end{equation}
centre relaxations, and therefore
\begin{equation}
    \Work\ \ge\ m\,N_{c}
    \ \ge\ \frac{3}{256\,\epsppr\sqrt{\alpha(1+\alpha)}},
    \label{eq:ladder-signed-work}
\end{equation}
with coefficient $3/128$ on the dyadic grid where $m=1/(8\epsppr)$ exactly.
\end{theorem}

\begin{proof}
$z_{c}^{0}=0$ and, by \eqref{eq:ladder-primitive}, $z_{c}$ changes only at
centre relaxations. The guarantee at $u=c$ forces
$z_{c}^{T}\ge\pi_{c}-\epsppr d_{c}=\pi_{c}-\epsppr m\ge\tfrac12-\tfrac18
=\tfrac38$, using \eqref{eq:ladder-star-closed-forms} and $\epsppr m\le1/8$.
The total travel $\sum\abs{\Delta z_{c}}=\sum\abs{\delta}$ over centre
operations is at least this, and each term is capped by
\Cref{lem:ladder-displacement}; each centre operation is charged
$d_{c}=m\ge1/(16\epsppr)$. The argument is pathwise, so randomized policies
satisfy it with probability one.
\end{proof}

\begin{theorem}[Matching member upper bound; \statusimported{} and
\statusmeasured]
\label{thm:ladder-signed-upper-bound}
On the same instance, SOR at $\omega^{\star}=1+\lambda^{2}$ with
$\lambda=(1-\sqrt\alpha)/(1+\sqrt\alpha)$, alternating centre and leaf blocks,
obeys the exact block law $x_{j}=(j+1)\lambda^{j}$ and stops, under the
residual certificate $\max_{u}\abs{r_{u}}/d_{u}\le\epsppr$ --- which implies
$\errop\le\epsppr$ by \ref{it:ladder-f4} and \ref{it:ladder-f5} --- within
\begin{equation}
    \Work_{\mathrm{SOR}}
    \ \le\ m\ceil{\frac{1}{\sqrt\alpha}
        \ln\frac{2}{\sqrt\alpha\,\epsppr\,m}}
    \ =\ \bigO\!\left(\frac{\log(1/\alpha)}{\sqrt\alpha\,\epsppr}\right).
    \label{eq:ladder-signed-upper}
\end{equation}
Consequently, on the star, the proved pair pins $\Work_{\Rpm}$ to
$\Theta\bigl(1/(\sqrt\alpha\,\epsppr)\bigr)$ up to the factor $\log(1/\alpha)$
carried by \eqref{eq:ladder-signed-upper}.
\end{theorem}

\noindent
The block law and \eqref{eq:ladder-signed-upper} are Lemma~3.1 and
Theorem~3.3 of \texttt{i2d} and are used without re-derivation. The
$\log(1/\alpha)$ factor is an artifact of the \emph{certificate} stop, not of
the method: under the semantic stop the measured block count satisfies
$N_{\mathrm{blocks}}\sqrt\alpha\in[1.35,1.50]$, flat over
$\alpha\in[2^{-12},2^{-4}]$, and the $\omega\to2$ alternating member attains
$\Work\sqrt\alpha\,\epsppr\in[0.083,0.094]$ for $\alpha\le2^{-5}$
(\statusmeasured; see \Cref{sec:ladder-verification}). Removing the logarithm therefore rests on \statusmeasured{} evidence, not on a
proof: with it, the two sides are pinned at ratio $3.5$ to $4.0$ and
$\Work_{\Rpm}=\Theta(1/(\sqrt\alpha\,\epsppr))$ exactly; without it, what is
proved is that same statement up to $\log(1/\alpha)$. The comparison with the
manuscript's certificate-stopped star bound is in
\Cref{sec:ladder-manuscript-relation}.

\begin{proposition}[Residual-mass cap, and output maps for $\Rpm$;
\statusproveddraft]
\label{prop:ladder-residual-mass-cap}
Along every $\Rpm$ path, on any graph and at every time,
\begin{equation}
    \norm{r}_{1}\ \le\ \sqrt{\vol(V)}\,\sqrt{2(1+\calpha)\Phi_{0}},
    \qquad\text{on the star}\qquad
    \norm{r}_{1}\ \le\ 2\sqrt{\frac{\alpha}{1+\alpha}} .
    \label{eq:ladder-residual-cap}
\end{equation}
Consequently \Cref{thm:ladder-signed-lower-bound} extends to outputs $z+Mr$
with $M\in\OB$: if $B\le\tfrac{1}{16}\sqrt{(1+\alpha)/\alpha}$ then the travel
is at least $1/4$ and
\begin{equation}
    N_{c}\ \ge\ \frac{1}{8\sqrt{\alpha(1+\alpha)}},
    \qquad
    \Work\ \ge\ \frac{m}{8\sqrt{\alpha(1+\alpha)}} .
    \label{eq:ladder-signed-outputmap}
\end{equation}
\end{proposition}

\begin{proof}
$D^{-1/2}r=\Hsym y$, so
$\norm{r}_{1}=\sum_{u}\sqrt{d_{u}}\,\abs{(\Hsym y)_{u}}
\le\sqrt{\vol(V)}\,\norm{\Hsym y}_{2}$ by Cauchy--Schwarz, and
$\norm{\Hsym y}_{2}^{2}=y^{\trans}\Hsym^{2}y\le\mu_{\max}y^{\trans}\Hsym y
=2\mu_{\max}\Phi\le2\mu_{\max}\Phi_{0}$ with $\mu_{\max}=1+\calpha$. The star
numbers follow from $\vol(V)=2m$ and $\Phi_{0}=\alpha/(2m)$. For the
extension, the guarantee at $c$ gives
$z_{c}^{T}\ge\pi_{c}-\epsppr m-(Mr)_{c}$ with
$(Mr)_{c}\le B\norm{r}_{1}\le2B\sqrt{\alpha/(1+\alpha)}\le1/8$ under the
stated bound on $B$, so the travel is at least
$\tfrac12-\tfrac18-\tfrac18=\tfrac14$; divide by
\eqref{eq:ladder-displacement-cap}.
\end{proof}

\begin{remark}[The admissible $B$ grows as $\alpha\to0$]
\label{rem:ladder-B-grows}
Energy forbids a signed-relaxation method from parking much mass in $r$, so
the tolerable column mass in \Cref{prop:ladder-residual-mass-cap} increases as
$\alpha$ decreases. Sharpness of that threshold is \emph{not} claimed; the
only upper delimiter asserted here is the universal zero-work escape of
\Cref{prop:ladder-escape}.
\end{remark}

\begin{proposition}[General graphs; \statusimported]
\label{prop:ladder-signed-general}
For any connected $G$ and seed $v$, let $P_{\ge\mu_{0}}$ project onto the
$\Hsym$-eigenspaces of eigenvalue at least $\mu_{0}$, and put
$\kappa_{v}:=\max_{\mu_{0}}\norm{P_{\ge\mu_{0}}\eunit{v}}\sqrt{\mu_{0}}$.
Then every seed relaxation satisfies
$\abs{\Delta z_{v}}\le2\sqrt{\galpha\pi_{v}}/\kappa_{v}$, so if
$\epsppr d_{v}\le\pi_{v}/2$ then
\begin{equation}
    N_{v}\ \ge\ \frac{\kappa_{v}}{4}\sqrt{\frac{\pi_{v}}{\galpha}},
    \qquad
    \Work\ \ge\ d_{v}N_{v}.
    \label{eq:ladder-signed-general}
\end{equation}
\end{proposition}

\noindent
This is Proposition~3.4 of \texttt{i4c}. The deciding instance parameter is
the seed's self-return amplification $\pi_{v}/\galpha$: it is
$\Theta(1/\alpha)$ on the star, where a degree-$\Theta(1/\epsppr)$ seed
reflects its mass back in one step, but only $\Theta(1/\sqrt\alpha)$ on paths,
spiders, caterpillars, trees and grids, where \eqref{eq:ladder-signed-general}
degrades to $\Omega(\alpha^{-1/4})$ seed operations (\statusmeasured, eleven
graphs). \Cref{open:ladder-beyond-star} asks whether this is intrinsic.
